\documentclass{beamer}
\usefonttheme{professionalfonts}
%\usetheme{Rochester}
%\usetheme{AnnArbor}
\usetheme{Szeged}
\useoutertheme{infolines}
\usecolortheme{dolphin}
\usepackage{graphicx}
% \usepackage[english]{babel}
\usepackage[dvipsnames]{xcolor} % opcional, para más colores
%\usepackage{hyperref}
\hypersetup{colorlinks=true, allcolors=blue}
%\definecolor{carmine}{rgb}{0.59, 0.0, 0.09}
\usepackage{amsmath} 

\usepackage{graphicx}
\usepackage{tikz}
\usetikzlibrary{shapes.callouts}
\usepackage{multirow}
\usepackage{booktabs}

\usepackage[utf8]{inputenc}

\usepackage{enumerate}
%\usepackage[shortlabels]{enumitem}
%\setitemize{itemsep=2pt, topsep=2pt}
\usepackage{amsmath,amsthm,amssymb,amsfonts,mathrsfs,latexsym,stmaryrd}
\usepackage{tcolorbox}

\usepackage{color} 
\usepackage{amssymb, amsmath, amsbsy}
\usepackage{adjustbox,lipsum}
\usepackage{placeins}

\usepackage{textpos}
\usepackage{wrapfig}
\usepackage[spanish, es-tabla]{babel}

\usepackage{caption}
\setlength{\parindent}{0pt}
%\usepackage[hyperfootnotes=false]{hyperref}

\usepackage{url}
\usepackage{float}
\usepackage{tikz}
\usepackage[table]{xcolor}
\usetikzlibrary{patterns}



\title[ICS-162]{Econometría ICS-294\\ \small{Clase 10: Inferencia sobre un parámetro}}

\author{Francisco Alfaro Medina}
\institute[USM]
{Universidad Técnica Federico Santa María\\ % Your institution for the title page
	\medskip
	\textit{} % Your email address
	
}
\logo{

\includegraphics[width=3cm]{sem7/images/Logo_UTFSM.png}
}
\date{}
\newcommand\Fontvi{\fontsize{9}{7.2}\selectfont}
\begin{document}
	
\begin{frame}
	\titlepage % Print the title page as the first slide


\end{frame}
%%%%%%%%%%%%%%%%%%%%%%%%%%%%%%%%%%%%%%%%%%%%%%%%%%%%%%%%%%%%%%%%%%%%%%%%%%%
\logo{}

% \begin{frame}
% \frametitle{Introducción}
% \textbf{Clase pasada:}
% \begin{itemize}
%     \item Planteamos los supuestos clásicos.
%     \item Mostramos que bajo los supuestos clásicos 1-5, el estimador MCO es insesgado.
%     \item Además, derivamos la forma matricial de la matriz de varianza-covarianzas y la varianza de cada estimador.
% \end{itemize}
% \textbf{Esta clase:}
% \begin{itemize}
%     \item Formula intuitiva para la varianza de los $\hat{\beta}$.
%     \item Función de distribución de los $\hat{\beta}$.
%     \item Test de hipótesis sobre un parámetro.
% \end{itemize}
% \end{frame}
% %%%%%%%%%%%%%%%%%%%%%%%%%%%%%%%%%%%%%%%%%%%%%%%%%%%%%%%%%%%%%%%%%%%%%%%%%%%
% \begin{frame}
% \frametitle{Varianza estimada de cada estimador $\beta_j$}
% Esta matriz de $(k + 1) \times (k + 1)$ se llama matriz de varianzas-covarianzas.
% \[
% \hat{Var}(\hat{\beta}|X) = \hat{\sigma}^2(X^{\top}X)^{-1}
% \]
% \begin{itemize}
%     \item El elemento de la diagonal $a_{1,1}$ es $\hat{Var}(\hat{\beta}_0)$.
%     \item El elemento de la diagonal $a_{2,2}$ es $\hat{Var}(\hat{\beta}_1)$.
%     \item Generalizando: el elemento de la diagonal $a_{j+1,j+1}$ es $\hat{Var}(\hat{\beta}_j)$.
% \end{itemize}
% \end{frame}

% \begin{frame}
% \frametitle{Ejemplo}
% Estimamos el siguiente modelo con dos variables:
% \[
% \text{salario}_i = \beta_0 + \beta_1 \text{educ}_i + \beta_2 \text{IQ}_i + u_i
% \]

%     \begin{figure}
%         \centering
%         \includegraphics[width=0.8\linewidth]{clase5/images/tabla0.png}
%     \end{figure}
    
% \end{frame}

% \begin{frame}[shrink=10]
% \frametitle{Varianza del Estimador MCO: Ejemplo}
% \begin{itemize}
%     \item Supuestos clásicos $\Rightarrow \hat{V}(\hat{\beta}) = \hat{\sigma}^2(X^{\top}X)^{-1}$.
%     \item Consideremos:
%     \[
%     (X^{\top}X)^{-1} =
%     \begin{bmatrix}
%     0.05987 & -0.00177 & -0.00035 \\
%     -0.00177 & 0.00030 & -0.00002 \\
%     -0.00035 & -0.00002 & 0.00001 \\
%     \end{bmatrix}\]
    
%      \[\hat{\sigma}^2 = \frac{SSR}{n-k-1} = 141925
%     \]
%     \item Entonces, podemos calcular:
%     \[
%     \hat{V}(\hat{\beta}) = \hat{\sigma}^2(X^{\top}X)^{-1} =
%     \begin{bmatrix}
%     8497.57 & -250.80 & -49.04 \\
%     -250.80 & 42.90 & -3.22 \\
%     -49.04 & -3.22 & 0.91 \\
%     \end{bmatrix}
%     \]
%     \item En la diagonal principal, tienen la varianza de cada uno de los estimadores.
%     \item Ejemplo: $\hat{V}(\hat{\beta}_1) = 42.9$. Es decir, los $\hat{se}(\hat{\beta}_1) = \sqrt{\hat{V}(\hat{\beta}_1)} = 6.54$.
% \end{itemize}
% \end{frame}

% \begin{frame}[shrink=10]
% \frametitle{Varianza del Estimador MCO}
% \begin{itemize}
%     \item Intuición de la varianza del $\hat{\beta}_1$ (elemento $a_{22}$ de la diagonal):
%     \[
%     \hat{Var}(\hat{\beta}_1|X) = \frac{\color{red}\hat{\sigma}^2}{SST_1(1 - R^2_1)}
%     \]
%     \item Donde:
%     \begin{itemize}
%         \item $SST_1 = \sum_{i=1}^{n}(x_{i1} - \bar{x}_1)^2 = n \cdot \text{Var}(x_1)$ es la suma total de los cuadrados de $x_1$.
%         \item $R^2_1$: el coeficiente de determinación resultante de la regresión de $x_1$ (como variable dependiente) en todas las otras variables independientes. Intuitivamente, indica cuánto de la variabilidad de $x_1$ se puede explicar por $x_2, x_3, \ldots$.
%     \end{itemize}
%     \item Ejemplo:
%     \begin{itemize}
%         \item Modelo: $y_i = \beta_0 + \beta_1x_{1i} + \beta_2x_{2i} + \beta_3x_{3i} + u_i$.
%         \item $R^2_1$ es el coeficiente de determinación de la regresión MCO auxiliar:
%         \[
%         x_{1i} = \alpha_0 + \alpha_2x_{2i} + \alpha_3x_{3i} + e_i
%         \]
%         $R^2_1$ indica cuánto de la variable $x_1$ es explicado por las otras variables independientes.
%     \end{itemize}
% \end{itemize}
% \end{frame}

% \begin{frame}[shrink=10]
% \frametitle{Varianza del estimador MCO: Modelo múltiple}
% Intuición de la varianza del $\hat{\beta}_1$:
% \[
% \text{Var}(\hat{\beta}_1|X) = \frac{\sigma^2}{\text{SST}_1(1 - R_{1}^{2})}
% \]
% \begin{itemize}
%     \item $\text{Var}(\hat{\beta}_1|X)$ es creciente en $\sigma^2$.
%     \begin{itemize}
%         \item Mientras más ruido no explicado haya en la ecuación, más difícil es estimar el efecto parcial de $x_1$ en $y$.
%     \end{itemize}
%     \item $\text{Var}(\hat{\beta}_1|X)$ es decreciente en $\text{SST}_1$.
%     \begin{itemize}
%         \item Esto implica que se prefiere una mayor variación muestral en $x_1$.
%         \item Una alternativa para aumentar esta cantidad y reducir la varianza es aumentar el tamaño muestral.
%         \item El supuesto 3 asegura que $\text{SST}_1 \neq 0$.
%     \end{itemize}
%     \item $\text{Var}(\hat{\beta}_1|X)$ es creciente en $R_{1}^{2}$.
%     \begin{itemize}
%         \item Intuición: Meter más variables le va quitando variabilidad a $x_1$. Esto hace que su estimador sea más impreciso.
%     \end{itemize}
% \end{itemize}
% \end{frame}

% \begin{frame}[shrink=10]
% \frametitle{Problema de Multicolinealidad}
% \[
% \text{Var}(\hat{\beta}_1|X) = \frac{\sigma^2}{SST_1(1 - \color{blue}R^2_1)}
% \]
% \begin{itemize}
%     \item Si dos variables son linealmente dependientes o están muy correlacionadas, la varianza del estimador será muy alta. En el límite, $R^2_1 = 1$ y la varianza es infinita. Esto viola el supuesto 3 de no multicolinealidad perfecta. En este caso, hay multicolinealidad perfecta.
%     \item Trade-offs:
%     \begin{enumerate}
%         \item \textbf{Varianza:} Incluir más variables explicativas \textcolor{red}{reduce la SSR ($\sigma^2$)}, pero \textcolor{blue}{aumenta $R^2_1$}. Efecto ambiguo sobre la varianza $\hat{\text{Var}}(\hat{\beta}_1)$.
%         \item \textbf{Sesgo:} Incluir más variables puede evitar el sesgo del coeficiente si la variable omitida correlaciona con la $x$. A menudo sacrificamos aumentar la varianza para estar seguros de evitar el sesgo.
%     \end{enumerate}
% \end{itemize}
% \end{frame}

\begin{frame}
\frametitle{Teorema de Gauss-Markov}
\begin{itemize}
    \item \textbf{Teorema:} Bajo los supuestos 1 a 5, el estimador de Mínimos Cuadrados Ordinarios (MCO) de $\beta$ es de mínima varianza entre los estimadores lineales insesgados de $\beta$.
    \item El estimador MCO es el estimador lineal insesgado de menor varianza (MELI / BLUE).
\end{itemize}
\end{frame}

\begin{frame}
\frametitle{Supuesto 6: Normalidad de $u$}
\begin{itemize}
    \item Las perturbaciones $u$ se distribuyen como $u \sim N(0, \sigma^2)$.
    \item Este supuesto es importante porque para realizar inferencia acerca de los $\beta$ y la construcción de intervalos de confianza se necesita la distribución de sus estimadores $\hat{\beta}$ (no basta con conocer su esperanza y varianza).
\end{itemize}
\end{frame}

\begin{frame}{Normalidad}
\begin{itemize}
    \item  Como u es la suma de muchos factores distintos no observados que afectan a y, se puede apelar
al teorema del límite central (TLC)  para concluir que u tiene una distribución
aproximadamente normal.

\item Se asume que todos los factores
no observados afectan a y por separado en forma aditiva.  Si u
es una  función no aditiva de los factores no observados, entonces el argumento del TLC no
puede emplearse.
\end{itemize}
\end{frame}
\begin{frame}{Supuestos del modelo lineal clásico (MLC)}

\begin{itemize}
\item Consisten en todos los supuestos de Gauss-Markov más el supuesto  del error
distribuido normalmente.
    \item Si se satisfacen los supuestos del MCL los estimadores de MCO son insesgados de varianza mínima, no es ya necesario restringir la
comparación a estimadores lineales en y.
\end{itemize}

{\color{blue} $$ y|X\sim N(\beta_0+\beta_1x_1+...+\beta_kx_k, \sigma^2)$$}
    
\end{frame}



\begin{frame}[shrink=10]
\frametitle{Distribución del estimador de MCO}
\begin{itemize}
    \item Teorema: bajo los supuestos 1 a 6 se cumple $\hat{\beta}_j|X \sim N(\beta_j, \text{Var}(\hat{\beta}_j|X))$.
    \item Esto implica que, condicional en las variables regresoras $X$,
    $$\frac{\hat{\beta}_j - \beta_j}{\sqrt{\text{Var}(\hat{\beta}_j|X)}} \sim N(0, 1).$$
    \item Pero, recuerden que no conocemos $\text{Var}(\hat{\beta}_j)$ porque no conocemos $\sigma^2$.
    \item En su lugar, vamos a usar la estimación: $\hat{\text{Var}}(\hat{\beta}_j)$.
    \item Nos queda que:
    $$\color{red}t_{\hat{\beta_j}}=\frac{\hat{\beta}_j - \beta_j}{\text{s.e.}(\hat{\beta}_j)} \sim t_{n-k-1}$$ 
    donde $\text{s.e.}(\hat{\beta}_j) = \sqrt{\hat{\text{Var}}(\hat{\beta}_j|X)}$.
    \item \textcolor{red}{El término en rojo va a ser clave para hacer inferencia y construir intervalos de confianza.}
\end{itemize}
\end{frame}

\begin{frame}[shrink=10]
\frametitle{Test de hipótesis sobre un parámetro individual}
\begin{itemize}
    \item Tenemos el modelo poblacional
    $$Y_i = \beta_0 + \beta_1X_{i1} + \ldots + \beta_kX_{ik} + u_i$$
    \item Hasta aquí, aprendimos a encontrar el estimador $\hat{\beta}$ de los $\beta$.
    \item Además, aprendimos a estimar su varianza y su función de distribución.
    \item Ahora queremos usar todo esto para testear hipótesis acerca del valor de $\beta$.
    
\end{itemize}
\end{frame}

\begin{frame}[shrink=10]
\frametitle{Test de hipótesis sobre un parámetro individual}
\begin{itemize}
    \item Suponga que quiere testear si el coeficiente $\beta_j$ que acompaña a $x_j$ es significativamente distinto al valor $a_j$. Donde $a_j$ es una constante (un número). Es decir:
    \begin{align*}
        H_0 &: \beta_j = a_j \\
        H_1 &: \beta_j \neq a_j
    \end{align*}
    \item {\color{red} Para el caso particular de $a_j = 0$, el test busca determinar si la variable $x_j$ es significativa o relevante para el modelo.

     Si se rechaza $H_0$, entonces $x_j$ es estadísticamente significativa.}
        
    \item Para el test se usa el siguiente estadístico $t$ construido bajo la hipótesis nula:
    $$t_{\hat{\beta}_j} = \frac{\hat{\beta}_j - a_j}{s.e.(\hat{\beta}_j)} \sim t_{n-k-1}$$
    \item $t_{\hat{\beta}_j}$ mide en cuántas desviaciones estándar estimadas de $\hat{\beta}_j$, se aleja $\hat{\beta}_j$ de $a_j$. Valores de $t_{\hat{\beta}_j}$ alejados de $a_j$ son evidencia en contra de la hipótesis nula.
\end{itemize}
\end{frame}

\begin{frame}{Error tipo I y Error tipo II}

\begin{table}[h!]
\centering
%\caption{Errores en una prueba de hipótesis}
\begin{tabular}{@{}llcc@{}}
\toprule
\textbf{Realidad} & \textbf{Decisión} & \textbf{Resultado} & \textbf{Tipo de error} \\ \midrule
$H_0$ verdadera & No se rechaza $H_0$ & Decisión correcta & -- \\
$H_0$ verdadera & Se rechaza $H_0$ & Decisión incorrecta & {\color{red}Error tipo I ($\alpha$) }\\
$H_0$ falsa     & No se rechaza $H_0$ & Decisión incorrecta & {\color{red}Error tipo II ($\beta$) }\\
$H_0$ falsa     & Se rechaza $H_0$    & Decisión correcta & -- \\
\bottomrule
\end{tabular}
\end{table}

    
\end{frame}


\begin{frame}
\frametitle{Errores tipo I y II }
    \renewcommand{\arraystretch}{1.3} % más espacio entre filas
\setlength{\tabcolsep}{5pt} % reducir espacio entre columnas

\rowcolors{2}{blue!5}{white} % alterna color de filas

\begin{tabular}{|p{2.8cm}|p{4cm}|p{4cm}|}
\hline
\rowcolor{blue!20}
\textbf{Característica} & \textbf{Error Tipo I ($\alpha$)} & \textbf{Error Tipo II ($\beta$)} \\
\hline
Definición & Rechazar $H_0$ cuando en realidad es verdadera. & No rechazar $H_0$ cuando en realidad es falsa. \\
\hline
% Interpretación & \textbf{Falso positivo}: concluir que hay un efecto cuando no lo hay. & \textbf{Falso negativo}: concluir que no hay efecto cuando sí lo hay. \\
% \hline
Control & Se fija con el \textbf{nivel de significancia} ($\alpha$). & Se reduce aumentando $n$. El \textbf{poder} de un test es  ($1-\beta$). \\
\hline
Consecuencia & Concluir efectos inexistentes. & No detectar un efecto real. \\

\hline
\end{tabular}

\end{frame}
\begin{frame}
\frametitle{Test de hipótesis sobre un parámetro individual}
\begin{itemize}
    \item {\color{blue}Supongamos que se fija la probabilidad de cometer error tipo I en $\alpha$. Nivel de significancia
o probabilidad de rechazar $H_0$ cuando en realidad es verdadera.}
    \item {\color{red}Sea $t_{\alpha/2}$ el percentil $(1 - \frac{\alpha}{2})100$ de la distribución t-Student de $n - k - 1$ grados de libertad.}
    \item Si $\left|\frac{\hat{\beta}_j - a_j}{s.e.(\hat{\beta}_j)}\right| \geq t_{\alpha/2}$, se rechaza $H_0$.
\end{itemize}

    \begin{figure}
        \centering
        \includegraphics[width=0.7\linewidth]{sem7/images/testhip.png}
    \end{figure}
    
\end{frame}



\begin{frame}
\frametitle{Test de hipótesis sobre un parámetro individual}
\begin{itemize}
    \item El anterior corresponde a un test de dos colas.
    \item Un nivel de significancia del $5\%$ es el valor más común.
    \item Esto significa que de todas las muestras aleatorias que se pueden tomar, en un $5\%$ de ellas se rechaza equivocadamente la hipótesis nula.
    \item Si se rechaza $H_0$ a un nivel de significancia de $\alpha$, entonces también se rechaza para niveles de significancia más altos. Por ejemplo, si una variable es significativa al $1\%$, entonces también lo es al $5\%$, al $10\%$, $20\%$, etc.
    \item \textcolor{blue}{Recuerde que $t_n \rightarrow N(0, 1)$ cuando $n \rightarrow \infty$. }
    %En la práctica, cuando los grados de libertad del estadístico t son más de 120 se suele aproximar la distribución t-Student con la Normal y usar los valores críticos de esta distribución.}
\end{itemize}
\end{frame}



% \begin{frame}{Valor-p}
% \begin{itemize}
%     \item Otra forma de reportar el rechazo - o no - de la hipótesis nula.
%     \item El p-valor es el menor nivel de significancia al que se habría rechazado la hipótesis nula.
%     \item Nos dice la probabilidad de rechazar incorrectamente la hipótesis nula.
%     \item Valores-p pequeños son evidencia en contra de la hipótesis nula.
% \end{itemize}
% \end{frame}

% \begin{frame}[shrink=10]{p-valor en gráfico}
%     \begin{itemize}
%         \item Ejemplo \( \hat{t}_{\hat{\beta}_1} = 2.4 \).
%         \item Buscar la probabilidad de que los valores puedan ser más extremos que \( \hat{t}_{\hat{\beta}_1} \) estadístico.
%     \end{itemize}

%     \begin{figure}
%         \centering
%         \includegraphics[width=0.8\linewidth]{clase5/images/pval.png}
%     \end{figure}
    
% \end{frame}

% \begin{frame}[shrink=10]{p-valor en gráfico}
%     \begin{itemize}
%         \item Valor del estadístico \( \hat{t}_{\hat{\beta}_1} = 2.4 \).
%         \item Buscar la probabilidad de que los valores puedan ser más extremos que \( \hat{t}_{\hat{\beta}_1} \).
%     \end{itemize}

%     \begin{figure}
%         \centering
%         \includegraphics[width=0.8\linewidth]{clase5/images/pval2.png}
%     \end{figure}

%     $p - value=0.016$
% \end{frame}

% \begin{frame}[shrink=10]{Ejemplo}
%     \begin{itemize}

%     En Stata, se reporta el t-estadístico y también el valor p.

%     \[ y_i = \beta_0 + \beta_1 \text{educ} + \beta_2 \text{IQ} + u_i \] 
    
%     \begin{figure}
%         \centering
%         \includegraphics[width=0.8\linewidth]{clase5/images/tabla.png}
%     \end{figure}
    
% Vemos que $p-value = 0.00$ para $\beta_1$. Que significa? Podemos rechazar que $\beta_1 = 0$ con un nivel de significancia de 0\%. O sea, seguro!
%     \end{itemize}
% \end{frame}

\begin{frame}{Intervalos de Confianza para un Parámetro Individual}
    \begin{itemize}
        \item El intervalo de confianza al $(1-\alpha)100\%$ para $\beta_j$ es:
        \[
        \hat{\beta}_j - t_{\alpha/2} \text{s.e.}(\hat{\beta}_j) \leq \beta_j \leq \hat{\beta}_j + t_{\alpha/2} \text{s.e.}(\hat{\beta}_j)
        \]
        \item Esto nos indica que con un $(1 - \alpha)100\%$ de confianza, $\beta_j$ estará dentro de ese intervalo.
        \item Si $\beta_j$ no está en el intervalo, entonces se rechaza la hipótesis nula al $(\alpha 100\%)$ de significancia.
    \end{itemize}
\end{frame}



\begin{frame}{Ejemplo: Datos}
Consideremos el modelo de regresión lineal:
\[
y_i = \beta_0 + \beta_1 \text{educ}_i + \beta_2 \text{IQ}_i + u_i
\]
donde queremos evaluar la hipótesis de que, después de tener en cuenta el IQ, la educación no tiene efecto en el salario.
\begin{figure}
    \centering
    \includegraphics[width=0.7\linewidth]{sem7/images/tabla.png}
\end{figure}

  Número de observaciones=935 
    
\end{frame}

\begin{frame}
\frametitle{Ejemplo de Test de Hipótesis}


\begin{enumerate}
    \item \textbf{Paso 1:} Planteamos las hipótesis:
    \[
    H_0 : \beta_1 = 0 \quad \text{vs} \quad H_1 : \beta_1 \neq 0
    \]
    
    \item \textbf{Paso 2:} Calculamos el estadístico t:
    \[
    t_{\hat{\beta}_1} = \frac{42.0576 - 0}{6.5498} = 6.4212
    \]
    
    \item \textbf{Paso 3:} Regla $\Rightarrow$ rechazamos $H_0$ si $|6.42| > t_{2.5}$, con gl=932.
    
    \item \textbf{Paso 4:} Dado que $|6.42| \geq 1.96$, rechazamos $H_0$. Con un nivel de significatividad del $5\%$, concluimos que $\beta_1$ es significativamente diferente de cero.
\end{enumerate}
\end{frame}
\begin{frame}[shrink=10]{Ejemplo}

   
    \begin{itemize}
        \item Tomemos $\alpha 100 = 5\% \rightarrow t_{\alpha/2} =  1.96$
        \item Recordando la fórmula:
        \[
        \hat{\beta}_1 - t_{\alpha/2} \text{s.e.}(\hat{\beta}_1) < \beta_1 < \hat{\beta}_1 + t_{\alpha/2} \text{s.e.}(\hat{\beta}_1)
        \]
        \item Reemplazando:
        \[
        42.0576 - 1.96 \times 6.5498 < \beta_1 < 42.0576 + 1.96 \times 6.5498
        \]
        \item Entonces:
        \[
        29.22 < \beta_1 < 54.895
        \]
        \item Esto significa que con 95\% de confianza, $\beta_1$ está entre esos valores.
    \end{itemize}

 \href{https://colab.research.google.com/drive/1ofQHrMDZSBYgfbw7JEjXtWvevykkMNX_?usp=sharing\#scrollTo=9Fdx9GNkxGIF}{LINK ejercicio}.   
\end{frame}


\begin{frame}{Valor-p (p-value)}


    ¿Cuál es el menor nivel de
significancia al que se habría rechazado la hipótesis nula? 

{\color{blue}Este nivel se conoce como el valor-p
de la prueba.}

El valor-p es la probabilidad de observar un
estadístico t tan extremo como el que se encontró si la hipótesis nula es verdadera.

Esto significa
que valores-p pequeños son evidencia contra la hipótesis nula.



\end{frame}

\begin{frame}{Valor-p}

\begin{figure}
    \centering
    \includegraphics[width=0.8\linewidth]{sem6//images/output (2).png}
\end{figure}
    
\end{frame}
\begin{frame}
\frametitle{Tests de una cola}
\begin{itemize}
    \item Estos tests consideran  hipótesis del tipo:
    \begin{itemize}
        \item $H_0 : \beta_j = a_j$
        \item $H_1 : \beta_j > a_j$
    \end{itemize}

    \item Valores negativos del estadístico $t _{\hat{\beta}_j}$ no representan evidencia en contra de la hipótesis nula.
\end{itemize}

\begin{figure}
    \centering
    \includegraphics[width=0.5\linewidth]{sem7/images/image.png}
\end{figure}


\end{frame}

\begin{frame}
\frametitle{Ejemplo: ecuación de salarios}
Consideremos el modelo
\[
\log(\text{wage}) = \beta_0 + \beta_1 \text{educ} + \beta_2 \text{exper} + \beta_3 \text{tenure} + u
\]
donde \textbf{educ} son años de educación, \textbf{exper} son años de experiencia y \textbf{tenure} son años de permanencia en la firma.


\begin{figure}
    \centering
    \includegraphics[width=0.6\linewidth]{sem7/images/salarios.png}
\end{figure}
 
\end{frame}

\begin{frame}
\frametitle{Ejemplo: ecuación de salarios}
Se busca testear si hay retornos positivos a la experiencia.

\begin{itemize}
    \item Paso 1: $H_0: \beta_2 =0$, $H_1: \beta_2 > 0$
    \item Paso 2: El estadístico $t$ es $t_{\hat{\beta}_2} = \frac{0.0041}{0.0017} \approx 2.41$
    \item Paso 3: Comparar $t_{\hat{\beta}_2}$ contra $t_{0.05}$. En este caso hay $n - k - 1 = 522$ gl, por lo que podemos usar valores de  la normal. El valor crítico a un nivel de significancia del 5\% es 1.645.
    \item Paso 4: Por lo tanto, $2.41 > 1.645$ y se concluye que $\beta_2$ es estadísticamente mayor que 0 con un nivel de significancia del 5\%.
\end{itemize}
\end{frame}

\begin{frame}
\frametitle{Ejercicio de cierre}

En el modelo de salarios se obtiene:
\[
\hat{\beta}_{educ}=0.075, \qquad se(\hat{\beta}_{educ})=0.025.
\]
Queremos testear \(H_0:\beta_{educ}=0\) contra \(H_1:\beta_{educ}\neq 0\), con \(\alpha=5\%\).

\begin{enumerate}
    \item Calcule el estadístico \(t\).
    \item Use el valor crítico aproximado \(1.96\). ¿Rechaza \(H_0\)?
    \item Interprete el resultado.
\end{enumerate}

\vspace{0.2cm}
\textbf{Respuesta:}
\[
t=\frac{0.075-0}{0.025}=3.
\]
Como \(|3|>1.96\), rechazamos \(H_0\). La educación tiene un efecto estadísticamente distinto de cero al 5\%.

\end{frame}

\begin{frame}
\frametitle{Conclusiones}
\begin{itemize}
    \item Hemos visto cómo se puede hacer tests de hipótesis simples sobre estimadores del MCO.
    \item El test más importante es determinar si el parámetro es significativamente diferente de cero.
    \item El estadístico \(t\) compara la distancia entre el estimador y el valor nulo en unidades de error estándar.
    \item Los intervalos de confianza y los valores-p entregan formas complementarias de reportar evidencia estadística.
   
\end{itemize}
\end{frame}


\end{document}
